\input ../preamble

\def\printmcount{\the\counter{section}.\the\counter{math counter}}
\setcounter{section}{2}

\centerline{\setfontandscale{bf}{20pt}Algebraic Geometry}
\medskip
\centerline{\bf Exercise \the\counter{section}}
\smallskip
\centerline{\it Ari Feiglin}

\bprob[title={Hartshorne II.2.4}]

    For a ring $A$ and a scheme $(X,\O_X)$ show that global sections is bijection
    $$ \hom_\Sch(X,\spec A) \xtos\Gamma \hom_\Ring(A,\Gamma(X,\O_X)) . $$

\eprob

We don't need that $X$ is a scheme, just that it is a locally ringed space.
For a morphism $\phi\colon X\to\spec A$ with underlying morphism of sheaves
$\phi^\sharp\colon\spec A\to\phi_*X$, for $f\in A$ let $\phi^\sharp_f$ be its $\dof(f)$-component.
So global sections maps $\phi$ to $\phi^\sharp_1$.
We want to show that $\phi^\sharp_1$ uniquely determines $\phi$; i.e.\ global sections is injective.

Let $x\in X$, then because $\phi$ is a local map we have that $\phi^\sharp$ is a local map of rings
$\O_{\spec A,\phi x}\to \O_{X,x}$.
So it maps the maximal ideal on the left into the maximal ideal of the right side; meaning
$\phi(x)A_{\phi x}$ is mapped into $\m_{X,x}$.
Now this means that for $f\in A$, $f\notin\phi(x)$ iff $f_x\notin\phi(x)A_{\phi x}$ iff it is a unit in
$\phi(x)A_{\phi x}$ iff $\phi_x(f_{\phi x})\notin\m_{X,x}$.
Here $f_{\phi x}$ and $\phi_x$ are the stalks of $f,\phi$ at $\phi x$ and $x$ respectively.

Because $\phi$ is a morphism, we have the following commutative diagram
\commsq{\Gamma(\spec A,\O_{\spec A})&\Gamma(X,\O_X)\cr \O_{\spec A,\phi x}&\O_{X,x}}
    {\phi^\sharp_1}{\phi^\sharp_x}{\bul_{\phi x}}{\bul_x}
Meaning $\phi^\sharp_x(f_{\phi x})=\phi^\sharp_1(f)_x$.
Thus
$$ f\in\phi(x) \iff \phi^\sharp_1(f)_x \in \m_{X,x} . $$
And so $\phi^\sharp_1$ uniquely determines the topological map $\phi$.

To show that it determines all of the other components of $\phi^\sharp$, it is sufficient to consider
only its components on distinguished open subsets $\dof f$, that is $\phi^\sharp_f$.
We have a commutative square (recall $\Gamma(\dof f,\spec A)=A_f$)
\commsq{A&\Gamma(X,\O_X)\cr A_f&\Gamma(\phi^{-1}\dof f,\O_X)}
    {\phi^\sharp_1}{\phi^\sharp_f}{\res}{\res}
By the universal property of localization, the bottom arrow is unique.
Thus $\phi^\sharp_1$ determines $\phi^\sharp_f$ as desired.
Note that in this proof we needed the fact that $\phi^\sharp_1$ also determines the topological map $\phi$
because we had to consider the pullback of $\O_X$.

So we have shown that global sections is injective.
It remains to show that it is surjective.

We follow closely the proof of injectiveness.
Suppose $\sigma\colon A\to\Gamma(X,\O_X)$ is a ring map.
First we define the topological map $\phi\colon X\to\spec A$, which we define according to the above proof.
For $x\in X$ let $\sigma_x\colon A\to\Gamma(X,\O_X)\to\O_{X,x}$ be the composition, which is a ring map,
and define $\phi(x)=\sigma_x^{-1}(\m_{X,x})$.
This is well-defined because $\sigma_x$ is a ring map and the preimage of a prime ideal is prime.

Now we claim $\sigma$ is continuous.
Indeed, for $f\in A$
$$ \sigma^{-1}\dof f = \set{x\in X}[f\notin\sigma_x^{-1}(\m_{X,x})] =
\set{x\in X}[\sigma_x(f)\notin\m_{X,x}] . $$
Note that $\sigma_x(f)$ is equal to the stalk of $\sigma(f)\in\Gamma(X,\O_X)$ at $x$.
So $\sigma^{-1}\dof f$ is the set of points $x\in X$ where the global section $\sigma(f)$ is nonzero
in the residue field $\kappa(x)=\O_{X,x}/\m_{X,x}$.
That is
$$ \sigma^{-1}\dof f = \set{x\in X}[\sigma(f)(x)\neq0] , $$
where $s(x)$ is equal to the image of $s_x\in\O_{X,x}$ under the projection $\O_{X,x}\to\kappa(x)$.

This is an open set: if $s\in\Gamma(X,\O_X)$ is nonzero in $\kappa(x)$, then it is invertible in
$\O_{X,x}$ (because it is local), so let $t\in\Gamma(U,\O_X)$ be a section whose germ is an inverse for
$s_x$.
Then we have $st=1$ in $\O_{X,x}$ which means that $st=1$ in a neighborhood of $x$, and thus $s$ is
invertible in a neighborhood of $x$, so for every $y$ in that neighborhood $s_y$ is invertible and thus
nonzero in $\kappa(y)$.
So this neighborhood is contained in $\set{x\in X}[s(x)\neq0]$, so it is open.

Thus $\sigma$ is continuous.
Now we define $\phi^\sharp\colon\O_{\spec A}\to\phi_*\O_X$ on the distinguished open sets.
We must of course have $\phi^\sharp_1=\sigma$, and by considering the commutative square (same as above)
\diagram{
    A&\Gamma(X,\O_X)\cr
    A_f&\Gamma(\phi^{-1}\dof f,\O_X)
}{
    \diagarrow{from={1,1}, to={1,2}, text={\phi^\sharp_1=\sigma}, y distance=.15cm}
    \diagarrow{from={2,1}, to={2,2}, text=\phi^\sharp_f, y distance=.15cm, dashed}
    \diagarrow{from={1,1}, to={2,1}, text=\res, x distance=.15cm}
    \diagarrow{from={1,2}, to={2,2}, text=\res, x distance=.15cm}
}
and by the universal property of localization, $\phi^\sharp_f$ is uniquely determined.
This forms a morphism of sheaves, again by the universal property.
Consider the diagram
\diagram{
    A&\Gamma(X,\O_X)\cr
    A_f&\Gamma(\phi^{-1}\dof f,\O_X)\cr
    A_g&\Gamma(\phi^{-1}\dof g,\O_X)\cr
}{
    \diagarrow{from={1,1}, to={1,2}, text={\phi^\sharp_1=\sigma}, y distance=.15cm}
    \diagarrow{from={2,1}, to={2,2}, text=\phi^\sharp_f, y distance=.15cm, dashed}
    \diagarrow{from={3,1}, to={3,2}, text=\phi^\sharp_g, y distance=.25cm, y off=.05cm, dashed}
    \diagarrow{from={3,1}, to={3,2}, text=\hat\phi^\sharp_g, y distance=-.3cm, y off=-.05cm, dashed}
    \diagarrow{from={1,1}, to={2,1}, text=\res, x distance=.15cm}
    \diagarrow{from={1,2}, to={2,2}, text=\res, x distance=.15cm}
    \diagarrow{from={2,1}, to={3,1}, text=\res, x distance=.15cm}
    \diagarrow{from={2,2}, to={3,2}, text=\res, x distance=.15cm}
}
Here $\phi^\sharp_1$ uniquely determines $\phi^\sharp_g$ and $\phi^\sharp_f$ determines a $\hat\phi^\sharp_g$.
So both squares here commute (the upper and lower).
These are again by localization.
But $\hat\phi^\sharp_g$ satisfies the universal property given by localization at $g$ for $A$ as well,
and thus it must coincide with $\phi^\sharp_g$, and so the middle horizontal arrows are equal and the
whole diagram commutes.

Thus this is a morphism of sheaves whose global section is $\sigma$.
We must show it is a local morphism, and then we will finally (finally) be done.
For $x\in X$ we need to show that $\phi(x)A_{\phi x}\leq\O_{\spec A,\phi x}$ is mapped into
$\m_{X,x}\leq\O_{X,x}$ by $\phi_x$.
Consider the commutative square
\commsq{A&\Gamma(X,\O_X)\cr A_{\phi x}&\O_{X,x}}
    {\sigma}{\phi_x}{\bul_{\phi x}}{\bul_x}
Then $\phi(x)A_{\phi x}$ is the image of $\phi(x)\leq A$ in $A_{\phi x}$ and so
$$ \phi_x(\phi(x)A_{\phi x}) = \sigma(\phi(x))_x = \sigma_x(\phi(x)) $$
(recall that $\sigma_x$ is the composition of $\sigma$ with taking the stalk at $x$).
But the definition of $\phi(x)$ is
$$ \phi(x) = \sigma_x^{-1}\m_{X,x} , $$
and thus
$$ \sigma_x(\phi(x)) \subseteq \m_{X,x} $$
as desired.
\qed

\bprob[title={Hartshorne II.2.14}]

    Let $S$ be a graded ring.
    \benum
        \item Show that $\Proj S=\emptyset$ iff every element of $S_+$ is nilpotent.
        \item Let $\phi\colon S\to T$ be a morphism of graded rings and define
        $U=\set{\p\in\Proj T}[\phi(S_+)\nsubseteq\p]$.
        Show that $U$ is open in $\Proj T$ and $\phi$ determines a natural morphism $\Phi\colon U\to\Proj S$.
        \item $\Phi$ may an isomorphism even when $\phi$ is not.
        Suppose that $\phi_d\colon S_d\to T_d$ is an isomorphism when $d\geq d_0$ for some integer $d_0$.
        Show that $U=\Proj T$ and $\Phi$ is an isomorphism.
    \eenum

\eprob

\benum
    \item If every element of $S_+$ is nilpotent, meaning $S_+\subseteq\nil(S)$, then every prime ideal of $S$ contains
    $S_+$ (because it contains $\nil(S)$).
    Thus every homogeneous prime ideal is irrelevant, and so $\Proj S=\emptyset$.

    Now we claim that
    $$ \nil(S) = \bigcap_{\mltext{$\p$ homogeneous\cr prime}}\p . $$
    Because $\nil(S)$ is the intersection of all prime ideals, it is contained in the intersection on the right.
    For an ideal $I$ define its {\emphcolor homogenization} to be the ideal generated by $I$'s homogeneous elements:
    $$ I_\homo = \ideal{a\in I}[\text{$a$ homogeneous}] = \bigoplus_n(I\cap S_n) . $$
    This is clearly a homogeneous ideal contained in $I$, and thus the homogeneous elements contained in $I_\homo$ are the same homogeneous elements
    contained in $I$.
    Also if $I$ is homogeneous then $I_\homo=I$.

    Furthermore for $\p$ prime, $\p_\homo$ is prime.
    This is because if $a,b\in S_\homo$ are homogeneous such that $a,b\notin\p_\homo$ then $a,b\notin\p$ because $\p$ and $\p_\homo$ have the
    same homogeneous elements, and thus $ab\notin\p$ so $ab\notin\p_\homo$.
    Because for homogeneous ideals primality can be checked on homogeneous elements, it follows that $\p_\homo$ is prime.

    Thus
    $$ \nil(S) = \bigcap_{\p\text{ prime}}\p \supseteq \bigcap_{\p\text{ prime}}\p_\homo = \bigcap_{\mltext{$\p$ homogeneous\cr prime}}\p , $$
    so we have equality.

    Now this means that if $\Proj S=\emptyset$ then all of $S$'s homogeneous prime ideals are irrelevant.
    Meaning for $\p$ homogeneous prime, $S_+\subseteq\p$.
    And so
    $$ S_+ \subseteq \bigcap_{\mltext{$\p$ homogeneous\cr prime}}\p = \nil(S) $$
    as desired.

    \item For a homogeneous prime $\p\leq T$, $\phi(S_+)\subseteq\p$ is equivalent to $\phi(f)\in\p$ for every homogeneous $f\in S_\homo$ of positive degree.
    Thus $\p\in U$ iff there is a positive degree $f\in S_\homo$ such that $\phi(f)\notin\p$, meaning
    $$ U = \bigcap_{f\in S_\homo}\dof_+(\phi f) , $$
    which is open.

    Now define $\Phi\colon U\to\Proj S$ by $\Phi(\p)=\phi^{-1}\p$.
    This is a continuous map, as for $f\in S_\homo$ of positive degree
    $$ \Phi^{-1}\dof_+(f) = \set{\p\in U}[f\notin\phi^{-1}\p] = \set{\p\in U}[\phi(f)\notin\p] = \dof_+(\phi f)\cap U , $$
    which is open.
    Note that $\dof_+(\phi f)\subseteq U$ so this is also just equal to $\dof_+(\phi f)$.

    We also define a morphism of sheaves $\Phi^\sharp\colon\O_{\Proj S}\to\Phi_*\O_U$ on the distinguished open sets of $\Proj S$.
    Its component on $\dof_+(f)\subseteq\Proj S$ is the map $\O_{\Proj S}(\dof_+(f))=S[1/f]_0$ to $\O_U(\dof_+(\phi f))=T[1/\phi f]_0$
    determined by $\phi$.
    In more detail, consider the composition $S\xto\phi T\to T[1/\phi f]$, and by localization this defines a map $S[1/f]\to T[1/\phi f]$.
    This map is graded because $\phi$ is ($a/f$ is sent to $\phi a/\phi f$, which has degree $\deg\phi a-\deg\phi f=\deg a-\deg f=\deg(a/f)$)
    and thus determines a morphism $S[1/f]_0\to T[1/\phi f]_0$.

    By universality of localization this is clearly a morphism of sheaves.
    That is, the restrictions $S[1/f]\to S[1/g]$ determine maps $S[1/g]\to T[1/\phi g]$ which by localization must be the same as the one
    determined by $S\to T$.

    \item First we show $U=\Proj T$.
    For $\p\leq T$ homogeneous prime, if $\phi(S_+)\subseteq\p$ then for $d\geq d_0$ we have $T_d\subseteq\p$.
    Now for $f\in T_d$ homogeneous of positive degree, we have that $f^{d_0}\in T_{dd_0}\subseteq\p$ and thus $f\in\p$.
    So $T_+\subseteq\p$, meaning $\p$ is irrelevant.
    Thus $\phi(S_+)\subseteq\p$ iff $T_+\subseteq\p$, which shows $U=\Proj T$ as desired.

    We also see that $\Phi$ is an open map onto its image as for $f\in T$ homogeneous of positive degree
    $$ \Phi\dof_+(f) = \set{\Phi\p}[f\notin\p] = \set{\Phi\p}[f^{d_0}\notin\p] \set{\Phi\p}[\phi^{-1}f^{d_0}\notin\Phi\p] = \dof_+(\phi^{-1}f^{d_0})\cap\im\Phi . $$

    Now $\Phi$ is injective as if $\phi^{-1}\q=\phi^{-1}\q'$ for $\q,\q'\in\Proj T$, then for $f\in T$ homogeneous of positive degree if $f\in\q$ then
    $f^{d_0}\in\q$ and there is an $x\in S$ such that $\phi(x)=f^{d_0}\in\q$ so that $\phi(x)\in\q'$ as well so $f^{d_0}\in\q'$ and thus $f\in\q'$.
    By symmetry $\q$ and $\q'$ have the same homogeneous elements of positive degree.
    The elements of positive degree determine the elements of zero degree, as $\q$ is relevant so if we choose $f\in T$ homogeneous of positive degree not in $\q$,
    then $g\in T_0$ is in $\q$ iff $fg\in\q$, and this is homogeneous of positive degree.
    Thus $\q$ and $\q'$ have the same homogeneous elements and are thus equal.

    And $\Phi$ is also surjective.
    For $\p\in\Proj S$ let us $f\in S$ be homogeneous of degree $\geq d_0$ such that $f\notin\p$.
    Now define
    $$ \q = \set{a\in T}[a\phi(f)\in\phi(\p)] . $$
    This is an ideal: it is clearly closed under addition and if $a\phi(f)\in\phi(\p)$ and $b\in T$ then $b\phi(f)$ has degree $\geq d_0$ and so lies in $S$.
    Thus $ab\phi f^2\in\phi(\p)$ (as it is an ideal), and since $f\notin\p$ by primality we have $ab\phi f\in\phi\p$ as desired.
    It is homogeneous because if $a\phi f\in\phi\p$ then taking $a_i\phi f=(a\phi f)_{i+d_0}$ and this is in $\phi\p$.
    Finally it is prime because if $a\phi f,b\phi f\notin\phi\p$ then $ab\phi f^2\notin\phi\p$ and thus $ab\phi f\notin\phi\p$.

    And $\phi(\p)\subseteq\q$ clearly so that $\p\subseteq\phi^{-1}\q$.
    Conversely if $\phi(a)\in\q$ then $\phi(af)\in\phi\p$ so $af\in\p$ and thus $a\in\p$.
    So $\phi^{-1}\q=\p$ as desired.
    Thus $\Phi$ is a bijective open continuous map, meaning it is a homeomorphism.

    And for $f$ of positive homogeneous, the natural map $S[1/f]\to T[1/\phi f]$ is an isomorphism.
    This is because elements on both sides may be written as $a/f^n$ and $b/\phi f^n$ for $a,b,f^n$ of degrees $\geq d_0$ (by exponentiating both the numerator and denominator
    by a sufficiently large $n$), and $\phi$ is an isomorphism for such degrees.
    And so $\phi^\sharp$ is an isomorphism on the distinguished open sets, and thus an isomorphism of sheaves as desired.
\eenum

\bprob[title={Hartshorne II.2.9}]

    A {\emphcolor generic point} for an irreducible space $Z$ is a point $p\in Z$ whose closure is all of $Z$.
    Show that for a closed irreducible subset $Z$ of a scheme $X$, $Z$ has a unique generic point.
    Such a space $X$ is called {\emphcolor Sober}.

\eprob

This is true for affine schemes, as we showed that closed irreducible subsets of $\spec A$ are the vanishing sets of prime ideals.
And the closure of $\p\in\spec A$ is precisely its vanishing set.
Furthermore the vanishing set of distinct prime ideals are distinct, showing uniqueness.

Now we make the following topological claim: if $X=\bigcup_iU_i$ is an open cover of a space $X$ such that each $U_i$ is Sober, then
$X$ is Sober.
Let $Z\subseteq X$ be irreducible and closed, then there is an $i$ for which $Z\cap U_i$ is nonempty.
As it is an open subset of $Z$, it is irreducible, and it is closed in $U_i$.
Thus there is some $p\in Z\cap U_i$ such that $Z\cap U_i=\clos_{U_i}(p)$.
Because $U_i\cap Z\subseteq Z$ is open and thus dense, we have that
$$ Z = \clos_Z(U_i\cap Z) = \clos_Z(\clos_{U_i}(p)) = \clos_Z(p) , $$
and because $Z$ is closed this means $Z=\clos_X(p)$, showing existence.

For uniqueness, if $Z=\clos_X(p)=\clos_X(q)$, then suppose $p\in U_i$.
If $q\notin U_i$ then $X-U_i$ is a closed set containing $q$ and so $\clos_X(q)\subseteq X-U_i$ which does not contain $p$ a contradiction.
Thus $p,q\in U_i$ and $\clos_{U_i}(p)=\clos_{U_i}(q)=Z\cap U_i$, and because $U_i$ is Sober this means $p=q$.

\bprob[title={Hartshorne II.3.1}]

    Show that a morphism of schemes $f\colon X\to Y$ is locally of finite type iff for {\it every\/} open affine $\spec B\subseteq Y$, $f^{-1}\spec B$
    can be covered with open affine $\spec A_i$ such that each $A_i$ is a finitely generated $B$-algebra.

\eprob

We appeal here to the affine communication lemma.
In particular we show that if $\spec B\subseteq Y$ is open affine satisfying the condition in question, then so is each $\spec B_x\subseteq\spec B$,
and if $(x_1,\dots,x_n)=B$ and $\spec B_{x_i}\subseteq Y$ satisfy the condition then so does $\spec B\subseteq Y$.

But first we use the affine communication lemma to prove that if $f^{-1}\spec B$ is covered by open affine $\spec A_i$ where each $A_i$
is finitely generated as a $B$-algebra, then every affine open subset of $f^{-1}\spec B$ is a finitely generated $B$-algebra.
So suppose $\spec A\xtos f\spec B$ corresponds to the finitely generated $B\xtos{f^*}A$.
Then for $x\in A$, $A_x$ is a finitely generated $A$-algebra, we have that $B\to A\to A_x$ is finitely generated.
Because $\spec A_f\inj\spec A$ corresponds to localization, we see that $\spec A_x\to\spec A\to\spec B$ corresponds to this map, as desired.
Now suppose that $A=(x_1,\dots,x_n)$ and $\spec A_{x_i}\xtos f\spec B$ are finitely generated.
This map corresponds to $B\xtos{f^*}A\to A_{x_i}$, so we need to show that if $A_{x_i}$ are finitely generated $B$-algebras, then so too is $A$ as desired.
This last claim follows from the fact that $A$ embeds into $\prod_iA_{x_i}$ via the natural map, so if the $A_{x_i}$ are finitely generated
then $A$ must be too.

Now we turn to the actual problem at hand.

Suppose that $\spec A\subseteq f^{-1}\spec B$ such that $f\colon\spec A\to\spec B$ corresponds to a finitely generated map $B\to A$.
Then $f\colon\spec A\to\spec B$ corresponds to a ring map $f^*\colon B\to A$, and $f|_{\spec A}^{-1}\spec B_x=\spec A_{f^*x}$.
And we know that if $B\to A$ is finitely generated, so too is $B_x\to A_{f^*x}$ (because it is generated by the generators of $A$).
Thus $f\colon\spec A_{f^*x}\to\spec B_x$ is finitely generated, and the $\spec A_{f^*x}$ cover $f^{-1}\spec B$ as desired.
Taking an open affine cover $f^{-1}\spec B=\bigcup_i\spec A_i$, we then have that $f^{-1}\spec B_x=\bigcup_i\spec A_{i,f_i^*x}$ and we showed that
$A_{i,f_i^*x}$ are finitely generated $B_x$-algebras as desired.

Conversely suppose $(x_1,\dots,x_n)=B$, and the condition is satisfied for $\spec B_{x_i}\subseteq Y$.
If $\spec A\subseteq f^{-1}\spec B$ is open affine, then let $f^*\colon B\to A$ be the associated ring map.
As before $\spec A_{f^*x_i}\subseteq f^{-1}\spec B_{x_i}$, and because the condition is satisfied for $\spec B_{x_i}$
we showed above that this means that every open affine subset of $f^{-1}\spec B_{x_i}$ is finitely generated over $B_{x_i}$.
So $B_{x_i}\to A_{f^*x_i}$ is finitely generated.
Because $B\to B_{x_i}$ is finitely generated, we have that $B\to A_{f^*x_i}$ is finitely generated.
And the $fx_i$ generate the unit ideal of $A$, so as before this means $B\to A$ is finitely generated.
And so the condition holds for $\spec B$.

So by the affine communication lemma, we have that if the condition is true for some open affine cover, it is true for every open affine subset.
In fact we have shown something even stronger: if $\spec B\subseteq Y$ and $\spec A\subseteq X$ are open affine with $f(\spec A)\subseteq B$,
then $A$ is a finitely generated $B$-algebra.

\bprob[title={Hartshorne II.3.14}]

    Let $X$ be a scheme of finite type over a field.
    Show that the set of closed points $\abs X$ is dense in $X$.

\eprob

The scheme $X$ being of finite type over a field $k$ means that for each open affine $\spec A\subseteq X$, $A$ is a finitely generated $k$-algebra.
This follows from the previous problem.

If $X$ is then an affine scheme, $X=\spec A$, this follows from a form of Nullstellensatz.
The closed points of $X$ are the maximal ideals of $A$, whose associated ideal is then $\jac(A)$ ($A$'s Jacobson radical).
By Nullstellensatz as $A$ is a finitely generated $k$-algebra, $\jac(A)=\nil(A)$.
The closure of $\abs X$ is the vanishing set of its associated ideal, which is $\vof\nil(A)=\spec A$.
Thus $\abs X$ is dense in $X$ as desired.

Now we claim that if $X$ is a scheme of finite type over $k$, then a closed point of an open affine subscheme of $X$ is a closed point of all of $X$.
That is if $\spec A\subseteq X$ is affine open and $\m\leq A$ is maximal, then $x=[\m]\in X$ is closed.
It is sufficient to show that if $x\in\spec B\subseteq X$ for any other affine open, then $x$ corresponds to a maximal ideal of $B$ as well.
This is sufficient as then $\set x$ is closed in all affine open subschemes, which cover $X$.

So suppose that $x$ corresponds to the prime ideal $\p\leq B$.
As $\spec A\cap\spec B$ is open, there is a $f\in B$ such that $x\in\spec B_f\subseteq\spec A$.
Because $x$ is closed in $\spec A$, it is closed in $\spec B_f$.
Thus $\p_f\leq B_f$ maximal.
And by assumption $B$ and $B_f$ are finitely generated $k$-algebras, and so the $k$-algebra morphism $B\to B_f$ pulls back maximal ideals of $B_f$
to maximal ideals of $B$.
The preimage of $\p_f\leq B_f$ is $\p$, and so $\p\leq B$ is maximal meaning $x$ is closed in $\spec B$ as well.

It then follows that the closed points of $X$ are the maximal ideals of its open affine subschemes.
In each open affine subscheme, the closure of this set is the whole subscheme, and thus its closure in $X$ is all of $X$ as desired.

\bprob

    Let $\phi\colon Z\to X$ be a morphism of schemes.
    Show that the two following conditions are equivalent.
    \benum
        \item $\phi$ is a closed embedding as a topological map, and $\phi^\sharp\colon\O_X\to\phi_*\O_Z$ is a surjective map of sheaves.
        \item For an affine open $\spec A\subseteq X$, $\phi^{-1}\spec A=\spec B$ is affine and $\phi\colon\spec B\to\spec A$ corresponds to
        a surjective map of rings $\phi^*\colon A\to B$.
    \eenum
    If $\phi$ satisfies any of these equivalent conditions, it is called a {\emphcolor closed immersion}.

\eprob

First we show (2) implies (1).
Cover $X$ with affine open $X=\bigcup_iU_i$, with $U_i=\spec A_i$ so that $\phi^{-1}\spec A_i=\spec B_i=V_i$.
Then $\phi\colon V_i\to U_i$ is a closed embedding for each $U_i$ as $B_i\to A_i$ is surjective (this is a claim regarding affine schemes which is clear).
Now we have a claim about topological spaces: if $X=\bigcup_iU_i$ is an open cover for which $\phi|_{\phi^{-1}U_i}\colon \phi^{-1}U_i\to U_i$ is a closed embedding,
then $\phi$ is a closed embedding.
It is clear that it is injective: for two elements with the same image, their image lies in some $U_i$, and thus both elements lie in $\phi^{-1}U_i$ for which
$\phi$ is injective and thus the elements are equal.
And $\phi$'s image is closed since $\phi(Z)\cap U_i=\phi(\phi^{-1}U_i)\subseteq U_i$ is closed for all $i$, and as the $U_i$ form an open cover, $\phi(Z)\subseteq X$
is closed.
Finally $\phi$ is an open map onto its image as this is true for an open cover of $Z$.

Then $\phi^\sharp|_{U_i}\colon\O_X|_{U_i}\to(\phi_*\O_Z)|_{U_i}\cong\phi_*(\O_Z|_{\phi^{-1}U_i})$ is surjective, as it corresponds to the map $\spec B_i\to\spec A_i$ and $A_i\to B_i$
is surjective.
%We know from the first problem set that for $p=\phi(z)\in\phi(Z)$, $(\phi_*\O_Z)_p=\O_{Z,z}$, and otherwise the stalk is zero.
Thus taking stalks shows that at each $p\in X$, $\phi^\sharp_z\colon\O_{X,p}\to\phi_*(\O_Z)_p$ is surjective (as $p$ must lie in some $U_i$, and restricted to $U_i$, $\phi^\sharp$ is
surjective).
And thus $\phi^\sharp\colon\O_X\to\phi_*\O_Z$ is surjective, as desired.

Now suppose (1), we will show (2).
First we show that when $X=\spec A$ then $Z=\spec B$ is affine and $\phi\colon\spec B\to\spec A$ corresponds to a surjective map of rings $A\to B$.

Let $Z=\bigcup_iV_i$ be an affine open covering with $V_i=\spec B_i$ such that $\phi(V_i)=\phi(Z)\cap U_i$ for $U_i\subseteq\spec A$ open.
We may refine $U_i$ to an open cover of distinguished open $\dof_A(f_{ij})\subseteq U_i$, and we cover the open $X-\phi(Y)$ with $\dof_A(f_{0j})$; since $X$ is quasicompact
we may take only finitely many $f_{ij}$.
Now $\phi|_{V_i}\colon V_i\to X$ is a morphism of affine schemes corresponding to $\phi_i\colon A\to B_i$, and the preimage of $\dof_A(f_{ij})$ under $\phi|_{V_i}$
is $\spec B_{i,\phi_if_{ij}}$.
Because $\dof_A(f_{ij})\cap\phi(Z)\subseteq\phi(V_i)$, we have that $\phi^{-1}\dof_A(f_{ij})$ is affine, equal to $\spec B_{i,\phi_if_{ij}}$.

Because $\dof_A(f_{ij})$ cover $\spec A$, the $f_{ij}$ generate the unit ideal of $A$, and so $f_{ij}\in\Gamma(Z,\O_Z)$ generate the unit ideal of $Z$ (here
we view $f_{ij}$ as its image under $\phi^\sharp$, which is equal to $\phi_if_{ij}$ when restricted to $V_i$).
Because $\phi^{-1}\dof_A(f_{ij})$ is affine, we have that $Z$ is affine.

Define for a scheme $X$, $A=\Gamma(X,\O_X)$.
The last claim is based off of the following facts
\benum
    \item if we let $X_f=\set{p\in X}[f(p)\neq0]$ (i.e.\ the open set of points where the global section $f$ is not in $\m_{X,p}$),
    then for $\phi\colon Y\to X$ then $\phi^{-1}X_f=Y_{\phi^\sharp f}$.
    \item if $X$ has a finite open affine cover $U_i$ such that $U_i\cap U_j$ is quasicompact, then $\Gamma(X_f,\O_X)\cong A_f$.
    \item if $\phi\colon X\to Y$ is a morphism of schemes such that there is an open cover $Y=\bigcup_iU_i$ such that $\phi\colon\phi^{-1}U_i\to U_i$
    is an isomorphism, then $\phi$ is an isomorphism.
    \item if $X$ is a scheme with global sections $f_1,\dots,f_n\in A$ which generate the unit ideal and $X_{f_i}$ are affine,
    then $X$ is affine.
\eenum
The first point is easy enough, because for $p\in Y$ we have $p\in Y_{\phi^\sharp f}$ iff $\phi^\sharp f_p\in\m_{Y,p}$ iff $f_{\phi p}\in\m_{X,\phi p}$ iff $\phi p\in X_f$.

For the second point, note that $X$ is quasicompact as the finite union of quasicompact subspaces.
So if $g\in A$ restricts to zero on $X_f$, then for affine $\spec B\subseteq X$ we have $g$ is zero in $X_f\cap\spec B=\dof_B(f)$.
And so $g$ is zero in $B_f$, which means there is an $n>0$ such that $f^ng=0$ in $\spec B$.
Covering $X$ with finitely many affine opens, we see there is an $n$ such that $f^ng=0$ on $X$.

Now let $U_i=\spec B_i$.
If $g\in\Gamma(X_f,\O_X)$ then in $U_i\cap X_f=\dof_{B_i}(f)$ we have that $f^ng\in B_i$ for some $n>0$.
So there is an $n>0$ such that $f^ng\in B_i$ for all $i$, more precisely there are $b_i\in B_i$ such that $b_i|_{\dof_{B_i}f}=(f^ng)|_{\dof_{B_i}f}$.
And in $U_i\cap U_j$ we have that both $b_i$ and $b_j$ are equal in $X_f$ to $f^ng$, and thus $b_i-b_j$ is zero in $X_f\cap U_i\cap U_j$.
So by the previous paragraph because $U_i\cap U_j$ is quasicompact there is an $m$ such that $f^m(b_i-b_j)=0$ on $U_i\cap U_J$.
Doing this for all $i,j$, by increasing $n$ we may assume that the $b_i$s agree on the overlaps $U_i\cap U_j$.
Thus they define a global section of the form $f^ng$.

So we see that there is a natural inclusion $\Gamma(X_f,\O_X)\to A_f$.
And for $h/f^n\in A_f$ ($h\in A$) we have that $h/f^n$ is in each $\dof_{B_i}(f)$ and so it is in $\Gamma(X_f,\O_X)$ so this is an isomorphism, proving (2).

For the third point, $\phi$ is a homeomorphism (this is just a topological claim, whose proof is similar to before).
And $\phi^\sharp|_{U_i}\colon\O_Y\to(\phi_*\O_X)|_{U_i}$ is an isomorphism as it corresponds to $\phi^{-1}U_i\to U_i$.
By taking stalks we see that for every $p\in Y$, $\phi^\sharp_p$ is an isomorphism (as $p\in U_i$ for some $i$), and thus $\phi^\sharp$ is an isomorphism
as desired.

For the fourth point let $A=\Gamma(X,\O_X)$ and consider the natural map $\phi\colon X\to\spec A$ induced by the global sections-spec adjunction and $\id\colon A\to A$.
Then $\phi^{-1}\spec A_{f_i}=X_{f_i}$ is affine by assumptuon.
By the second point $\phi$ defines an isomorphism on the global sections, and because these are affine schemes, $\phi$ restricts to an isomorphism $\phi^{-1}\spec A_{f_i}\to X_{f_i}$.
The $X_{f_i}$ cover $X$ and thus by the third point $\phi$ is an isomorphism, so $X$ is affine as desired.

\bye

